% ECONOMICS_PROBLEM: {"id": "D71-190", "title": "The least universal size of a Condorcet-winning set", "jel_code": "D71", "source_id": "OP-0203"}
% !TeX program = xelatex
\documentclass[11pt,a4paper]{article}
\usepackage[margin=23mm]{geometry}
\usepackage{fontspec}
\setmainfont{Latin Modern Roman}
\usepackage{amsmath,amssymb,mathtools}
\usepackage{xcolor,enumitem,booktabs,longtable,array,xurl,hyperref}
\definecolor{muted}{HTML}{53636C}
\hypersetup{colorlinks=true,urlcolor=blue,linkcolor=blue}
\setlength{\parindent}{0pt}
\setlength{\parskip}{5pt}
\newcommand{\R}{\mathbb R}
\newcommand{\E}{\mathbb E}
\newcommand{\N}{\mathbb N}
\newcommand{\ind}{\mathbf 1}
\newcommand{\Var}{\operatorname{Var}}
\newcommand{\Cov}{\operatorname{Cov}}
\newcommand{\supp}{\operatorname{supp}}
\DeclareMathOperator*{\argmin}{arg\,min}
\DeclareMathOperator*{\argmax}{arg\,max}
\newcommand{\entrymeta}[3]{{\sffamily\small\color{muted}\textbf{\texttt{#1}}\quad|\quad #2\par #3\par}}
\newcommand{\Problemref}[1]{\texttt{#1}}
\newcommand{\JELref}[2]{\texttt{#2} (JEL \texttt{#1})}
\begin{document}
\section*{The least universal size of a Condorcet-winning set}
\textbf{Problem ID:} \texttt{D71-190}\quad \textbf{JEL:} \texttt{D71}\par
% BEGIN ECONOMICS STATEMENT
\label{problem:OP-0203}
% GAME_THEORY_PROBLEM: {"local_id": "GT-073", "original_number": 73, "kind": "P", "entry_kind": "statement_only", "published": false, "review_required": true, "category": "game_theory"}
\label{game:GT-073}
{\sffamily\small\color{muted}
\noindent\textbf{ID: GT-073.}\quad\textbf{Category:} Social choice.
\quad\textbf{Kind: P.}\quad\textbf{Status:} Open exact-value problem;
the catalogue's upper bound is corrected.
\par}
\subsection*{Definitions and mathematical statement}
An election consists of a finite nonempty voter set \(N\), a finite
nonempty candidate set \(C\), and a strict total order \(\succ_i\) on
\(C\) for every \(i\in N\). A nonempty set \(S\subseteq C\) is
\emph{Condorcet-winning} if
\[
 \left|\left\{i\in N:
          a\succ_i s\text{ for every }s\in S\right\}\right|
 \leq \frac{|N|}{2}
 \qquad(a\in C\setminus S).
\]
Thus no outside candidate is ranked above the entire set by a strict
majority. Define
\[
 d(N,C,\succ)=
 \min\{|S|:\varnothing\ne S\subseteq C,\ S\text{ is Condorcet-winning}\},
 \qquad
 k^\star=\sup_{\text{finite elections}}d(N,C,\succ).
\]
The minimum exists because \(S=C\) satisfies the condition.

\paragraph{Question.}
Determine the exact value of \(k^\star\). The bounds verified in [S1] are
\[
                         3\leq k^\star\leq 5.
\]

\subsection*{Short English statement}
What is the smallest fixed number of candidates sufficient to obtain
a Condorcet-winning set in every election?

\subsection*{Variants and scope boundaries}
The majority compares an outsider with every member of \(S\) using the
same voters. This does not require some single member of \(S\) to beat
each outsider in a pairwise majority contest.
An exact half of voters does not constitute a blocking strict majority.

\subsection*{Source relationship and status}
[S1, Theorem 1] improves the upper bound to five.
The catalogue incorrectly reported the previously known upper bound
of six as the paper's final bound.

\subsection*{References}
\begingroup\footnotesize\raggedright
\noindent[S1] Thanh Nguyen, Haoyu Song, and Young-San Lin,
\emph{A Few Good Choices}, 2025, arXiv:2506.22133v2.
\url{https://arxiv.org/pdf/2506.22133}.
Locators: Section 1, printed pp.~2--3, including Theorem 1;
related-work discussion, printed p.~5 (upper bound five and lower
bound three).
\par\endgroup

\subsection*{Common conventions: Source mathematical conventions}

\textbf{Finite games.} For a finite set $A$, $\Delta(A)$ is its probability
simplex. Players choose independent mixed strategies in Nash equilibrium;
correlation is allowed only when explicitly part of the solution concept.
In a game with payoff functions $u_i$ and strategy profile $x$, an additive
$\varepsilon$ Nash equilibrium satisfies
\[
 u_i(x)\geq u_i(x_i',x_{-i})-\varepsilon
 \quad\text{for every player $i$ and every admissible deviation $x_i'$}.
\]
For numerical approximation, payoffs are normalized as locally specified.
An $\varepsilon$ payoff residual is distinct from distance at most
$\varepsilon$ to the set of exact equilibria. Point convergence, convergence
of empirical play, and convergence in distance to an equilibrium set are
also distinct.

\textbf{Computational representations.} Unless stated otherwise, a complexity
question takes rational data with binary encoding; polynomial time is measured
in total bit length. A fixed-accuracy polynomial algorithm is not necessarily
polynomial in $1/\varepsilon$, and neither is necessarily polynomial in
$\log(1/\varepsilon)$. Succinct games and value, demand, or sampling oracles
must declare their interfaces. Exact real solutions require an explicit
representation; an unspecified list of arbitrary real numbers is not a
finite computational output. A claimed algorithmic barrier is meaningful
only for its specified model and reductions.

\textbf{Dynamic games.} Histories, information, observation, and allowable
randomization are part of the model. A stationary strategy depends only on
the current state; a positional strategy in a deterministic graph game
chooses one action at each controlled vertex. Nash equilibrium at an initial
state and equilibrium after every history are different requirements.
An expectation of a pathwise $\liminf$ payoff, a $\liminf$ of expected averages,
a limiting expected average, and a uniform finite-horizon equilibrium are
different criteria. None is silently substituted for another.

\textbf{Allocation and mechanisms.} A complete allocation assigns every item
exactly once unless otherwise stated. Zero-valued goods, negative values,
capacity constraints, feasibility, lotteries, and copies of goods affect
fairness definitions and are specified locally. Dominant-strategy incentive
compatibility, Bayesian incentive compatibility, universal truthfulness,
and truthfulness in expectation impose different inequalities. Whenever
an objective ratio has zero denominator, its convention is stated or those
instances are excluded.

\textbf{Infinite-dimensional models.} Expectations, integrals, stochastic
equations, and optimization objectives require their local regularity,
integrability, filtrations, and admissible domains. A backward--forward
mean-field system is a boundary-value problem; it is not an initial-value
semiflow without an additional construction. Long-time, large-population,
and vanishing-noise limits require an order of limits and a convergence
topology. If a minimum need not be attained, the objective is its infimum
or supremum and the approximation requirement is stated separately.
% END ECONOMICS STATEMENT
\end{document}
